\documentclass[11pt]{article}
\usepackage[margin=1in]{geometry}
\usepackage{booktabs}
\usepackage{hyperref}
\usepackage{amsmath}
\usepackage{amssymb}
\usepackage{url}
\usepackage{xcolor}
\usepackage{enumerate}
\usepackage{float}

\title{Memo 07: Minimum-Effort Sediment Floc Simulation\\
       Using the Existing FVCOM Sediment Capability}
\author{FVCOM-MP WaterPACT Investigation}
\date{May 2026}

\begin{document}
\sloppy
\maketitle

\section*{Purpose}

This memo records a code-level investigation of how to simulate sediment floc
behavior by itself in the WaterPACT estuary setup, with minimum code and input
disruption.  The user request referred to \texttt{sed\_mod.F}; in this
repository there is no file with that exact name.  The relevant sediment source
files are:
\begin{itemize}
  \item \texttt{mod\_sed.F}: the original FVCOM sediment module, enabled by
        \texttt{-DSEDIMENT -DORIG\_SED}.
  \item \texttt{mod\_sed\_cstms.F} and \texttt{mod\_flocmod.F}: the CSTMS
        sediment and floc population-balance modules, enabled by
        \texttt{-DSEDIMENT -DCSTMS\_SED}.
\end{itemize}

The practical question is whether we can run only sediment flocculation in the
estuary without activating the microplastic floc-interaction machinery and
without moving to a new sediment framework.

\section*{Main Finding}

The minimum-effort path is to keep the original FVCOM sediment module
(\texttt{ORIG\_SED}) and change the sediment class from
\texttt{non-cohesive} to \texttt{cohesive}.  In \texttt{mod\_sed.F}, the
cohesive branch of \texttt{Calc\_Deposition} automatically calls
\texttt{Calc\_Wset\_Floc}; the non-cohesive branch uses the constant
\texttt{SED\_WSET}.  Therefore, the current WaterPACT sediment files do not
exercise the sediment floc settling logic because they set:
\begin{verbatim}
SED_TYPE = non-cohesive
\end{verbatim}

Changing that one semantic switch, with a physically smaller
\texttt{SED\_SD50} and \texttt{SED\_WSET}, gives a sediment-only floc
experiment using existing code.

\section*{What the Existing ORIG Sediment Floc Path Does}

For each cohesive sediment class, \texttt{mod\_sed.F} computes a local settling
velocity multiplier from suspended sediment concentration, salinity, turbulence
intensity, and primary grain size:
\[
  W_s = W_{s0}\,
        \frac{a\,C^b\,S^d\,G^e}{(D_{50,\mathrm{mm}})^f},
\]
with default constants in the source:
\[
  a=0.274,\quad b=0.48,\quad d=0.03,\quad e=0.22,\quad f=0.58.
\]
The code enforces a minimum multiplier of 1.0 and only allows enhancement in
brackish water:
\[
  5 < S < 15.
\]
It also caps the settling velocity at \texttt{0.5e-3 m/s}, or
0.5 mm/s.  This matters because the current WaterPACT sediment files use
\texttt{SED\_WSET = 3.0} mm/s for sand-like particles; if that value is reused
for cohesive sediment, the result will immediately hit the cap.  A floc-only
cohesive test should use a smaller base settling speed, for example
\texttt{SED\_WSET = 0.05--0.2} mm/s.

This path is lightweight but limited: it represents flocculation as a
diagnostic settling-velocity adjustment, not as an explicit floc size
distribution.

\subsection*{Meaning of the Formula Variables}

In the source routine \texttt{Calc\_Wset\_Floc(n,c,s,g1,ds50,w0,wset)}, the
variables in the empirical multiplier are:
\begin{itemize}
  \item \texttt{c}: suspended sediment concentration in the water column.  In
        the ORIG sediment module this is stored internally as kg\,m$^{-3}$,
        which is numerically equal to g\,L$^{-1}$.
  \item \texttt{s}: salinity, from \texttt{s1(i,k)}.
  \item \texttt{g1}: a turbulence-intensity/shear proxy computed immediately
        before the call from local velocity, water depth, the Manning-like
        coefficient \texttt{mf=0.012}, gravity, and \texttt{Vcons=1.e-6}.  It
        is not a direct GOTM/GLS dissipation rate; it is the ORIG module's
        empirical turbulence-intensity argument for this formula.
  \item \texttt{ds50}: primary sediment diameter after conversion to meters.
        In the formula it is multiplied by 1000, so the denominator uses
        \texttt{D50} in mm.
  \item \texttt{w0}: base settling velocity without the floc multiplier.
  \item \texttt{wset}: final settling velocity passed into the vertical
        settling calculation.
\end{itemize}
So, yes: in the shorthand interpretation of the formula, \(C\) is concentration,
\(S\) is salinity, and \(G\) is the local turbulence/shear-intensity argument
used by this empirical ORIG sediment closure.

\section*{Can This Be Used To Back Out an Equivalent Floc Size?}

Only with additional assumptions.  The ORIG settling-velocity formulation does
not prognose floc diameter, floc density, porosity, or fractal dimension.  It
returns an effective settling velocity.  A unique floc size cannot be recovered
from settling velocity alone because floc settling depends on both diameter and
effective density, and floc effective density itself usually varies with floc
size and composition.

For post-processing, one could compute a diagnostic ``Stokes-equivalent''
diameter by assuming a spherical floc with a specified effective density:
\[
  d_{\mathrm{eq}} =
  \left(\frac{18\mu W_s}{g(\rho_f-\rho_w)}\right)^{1/2}.
\]
This is useful only as an interpretation aid.  It is sensitive to the assumed
\(\rho_f-\rho_w\), and it becomes especially ambiguous near the hard cap
\(W_s \le 0.5\) mm\,s$^{-1}$.

A more floc-aware diagnostic would assume a primary particle diameter
\(d_p\), sediment density \(\rho_s\), water density \(\rho_w\), and fractal
dimension \(n_f\), for example:
\[
  \rho_f-\rho_w = (\rho_s-\rho_w)
  \left(\frac{d_p}{d_f}\right)^{3-n_f}.
\]
Under Stokes-flow assumptions this gives:
\[
  d_f =
  \left[
  \frac{18\mu W_s}{g(\rho_s-\rho_w)d_p^{\,3-n_f}}
  \right]^{1/(n_f-1)}.
\]
That calculation would still be diagnostic, not part of the current model
state.  If explicit floc size is a primary output requirement, the CSTMS
\texttt{mod\_flocmod.F} route is the more appropriate model structure.

\section*{Planned MP--Floc Coupling Extension}

The preferred future coupling plan is to keep the simplified ORIG sediment
floc module internally unchanged, while using a diagnostic/effective floc size
only inside the plastic--floc interaction pathway.  In other words:
\begin{itemize}
  \item \texttt{SED\_SD50} remains the constant primary cohesive sediment size
        used by \texttt{mod\_sed.F}.
  \item The sediment floc formula continues to compute the local
        floc-affected settling velocity and stores it internally as
        \texttt{sed(ised)\%wset0(i,k)}.
  \item \texttt{plast(iplast)\%dequ} remains the true MP equivalent particle
        diameter.  It should not be overwritten by a floc diameter.
  \item A separate diagnostic variable, conceptually
        \texttt{d\_floc\_eff(i,k,ised)}, would represent the effective floc
        size used only for MP--floc aggregation, size-exclusion, and related
        plastic interaction terms.
\end{itemize}

This is a one-way coupling approximation: sediment concentration, salinity,
turbulence, and primary sediment size determine the floc environment; MP then
responds to that environment.  Because MP concentration is much smaller than
suspended sediment concentration in the intended application, MP does not need
to feed back on the sediment floc state in this minimum implementation.

\subsection*{Current Behavior to Preserve}

In the present \texttt{mod\_plast.F} floc-interaction pathway, the switch
\texttt{PLAST\_FLOC} activates the two-tracer aggregated/disaggregated MP
logic.  The current implementation uses sediment \texttt{D50}-based proxies:
\begin{itemize}
  \item the MP--floc aggregation kernel uses \texttt{sed(ised)\%Sd50} as the
        floc-size proxy;
  \item the differential-settling term uses the class settling velocity
        \texttt{sed(ised)\%Wset};
  \item aggregated MP settling currently uses the sediment settling velocity
        proxy already in the code.
\end{itemize}
This current behavior should remain exactly available as the baseline option.

\subsection*{Proposed New Switch}

Add one additional plastic-module runtime switch after the literature
parameterization is selected, for example:
\begin{verbatim}
PLAST_FLOC_EFFECTIVE = F
\end{verbatim}
The exact name can be chosen during implementation.  The required behavior is:
\begin{itemize}
  \item \textbf{When false:} preserve the current \texttt{mod\_plast.F}
        behavior exactly.  MP--floc interaction uses sediment
        \texttt{Sd50}/\texttt{Wset} proxies, with no diagnostic effective floc
        size.
  \item \textbf{When true:} use the modified floc quantities only inside the
        plastic--floc interaction pathway.  The aggregation and size-exclusion
        calculations should use \texttt{d\_floc\_eff(i,k,ised)}, and the
        aggregated MP settling calculation should use the floc-affected
        sediment settling velocity, preferably \texttt{sed(ised)\%wset0(i,k)}
        or a sediment-concentration-weighted equivalent if multiple cohesive
        classes are active.
\end{itemize}

The open science item is the parameterization for
\texttt{d\_floc\_eff}.  The simple Stokes or fractal inversion written above is
not yet satisfactory as a production closure.  Before implementing this switch,
we should choose a defensible literature-based relation for effective floc
diameter and, if needed, floc effective density/porosity.  That reference
choice should also define how sediment mass concentration is converted to floc
number density in the MP aggregation kernel.

\section*{Reference Trail for the ORIG Formula}

The immediate source-code reference is the comment block above
\texttt{Calc\_Wset\_Floc}.  It attributes the cohesive floc settling method to
Qingyu Sha and lists:
\begin{itemize}
  \item Ding P.X., Hu K.L., Kong Y.Z., Zhu S.X. et al. (2003).
        ``Numerical simulation of total sediment under waves and currents in
        the Changjiang Estuary.'' \textit{Acta Oceanologica Sinica},
        25(5):113--124.  The same citation is visible in online bibliographic
        lists and environmental reports as a Chinese article with an English
        abstract.
  \item Hu K.L. (2003). ``2-D Suspended Sediment Numerical Simulation under
        Waves and Currents in Yangtze Estuary.'' Ph.D. dissertation, East
        China Normal University, as cited by the source comment.
\end{itemize}

I could verify the bibliographic identity of the Ding et al. paper online, but
I did not find an openly accessible copy containing the exact empirical
coefficient equation.  The strongest citation for the exact coefficient set is
therefore the FVCOM source comment itself, which appears to have embedded the
Changjiang/Yangtze-estuary empirical relation directly from the cited Chinese
paper/dissertation trail.

Online sources checked include a later Yangtze Estuary process-modeling paper
whose reference list includes Ding et al. (2003), and a Changjiang Estuary
flocculation review/survey PDF.  Those sources support the same physical
dependencies used by the code: concentration, salinity, particle size, and
hydrodynamic/turbulence conditions control floc formation and settling.  Recent
summaries also place an important Changjiang flocculation salinity range around
5--10, which is consistent with the source code's brackish activation gate of
5--15.
\begin{itemize}
  \item Process-modeling reference list checked:
        \url{https://www.researchgate.net/publication/292835947_A_process-based_approach_to_sediment_transport_in_the_Yangtze_Estuary}
  \item Changjiang Estuary flocculation survey PDF checked:
        \url{https://zjjcmspublic.oss-cn-hangzhou-zwynet-d01-a.internet.cloud.zj.gov.cn/jcms_files/jcms1/web2548/site/attach/0/15fb3123b5794fea9144293da397749a.pdf}
\end{itemize}

\section*{Current WaterPACT Configuration Relevant to Floc-Only Runs}

\begin{table}[H]
\centering
\caption{Current input features that matter for a sediment-only floc test.}
\small
\begin{tabular}{p{0.20\textwidth}p{0.34\textwidth}p{0.36\textwidth}}
\toprule
Item & Current setting & Implication \\
\midrule
Run namelist & \texttt{SEDIMENT\_MODEL = T} & Sediment runtime switch is already used. \\
Run namelist & \texttt{PLASTIC\_MODEL = T} & Turn this off for sediment-only floc. \\
Sediment file & \texttt{SED\_TYPE = non-cohesive} & Floc settling path is inactive. \\
Sediment file & \texttt{SED\_WSET = 3.0} mm/s & Too large for cohesive floc branch. \\
Sediment file & \texttt{MORPHO\_FACTOR = 10.0} & Accelerates bed exchange; use 1.0 for cleaner physics. \\
Sediment file & \texttt{INIT\_BED\_FRACTION} has six entries & For \texttt{NSED=1}, set it to 1.0. \\
River file & variable \texttt{coarse\_sand\_1} exists & Keep this sediment name to avoid NetCDF edits. \\
\bottomrule
\end{tabular}
\end{table}

The existing river NetCDF file \texttt{waterPACT\_riv\_floc\_MP.nc} already
contains a sediment concentration variable named \texttt{coarse\_sand\_1},
with values about \texttt{4.5e-4} to \texttt{1.6e-2} g/L.  FVCOM looks up
river sediment variables by the sediment class names passed from
\texttt{SETUP\_SED}; therefore, the lowest-effort route is to keep
\texttt{SED\_NAME = coarse\_sand\_1} even if the class is reinterpreted as
cohesive mud/floc material.  If the name is changed to \texttt{mud\_floc} or
similar, the river NetCDF file must also contain a matching variable.

\section*{Recommended Minimum-Effort Setup}

\subsection*{1. Use the ORIG sediment build}

For a fresh compile, make sure the executable is built with the original
sediment module:
\begin{verbatim}
FLAG_21  = -DSEDIMENT
FLAG_211 = -DORIG_SED
\end{verbatim}
The checked \texttt{make.inc} currently has these sediment flags commented,
while \texttt{FLAG\_49 = -DPLASTIC} is enabled.  Existing executables may have
been compiled from a different flag state, so this should be confirmed before
rebuilding or submitting a new production run.

\subsection*{2. Turn off plastic in the run namelist}

Clone one of the current WaterPACT run directories and change only the model
switches needed to remove microplastic:
\begin{verbatim}
SEDIMENT_MODEL          = T,
SEDIMENT_MODEL_FILE     = 'generic_sediment_floc_only.inp',
SEDIMENT_PARAMETER_TYPE = 'uniform',
SEDIMENT_PARAMETER_FILE = 'none',
PLASTIC_MODEL           = F,
PLASTIC_MODEL_FILE      = 'none'
\end{verbatim}

Keep the existing river forcing:
\begin{verbatim}
RIVER_NUMBER       = 9,
RIVER_INFO_FILE    = 'waterPACT_riv_floc_MP.nml',
RIVER_KIND         = 'variable',
RIVER_TS_SETTING   = 'calculated',
RIVER_INFLOW_LOCATION = 'node'
\end{verbatim}
The river file contains both sediment and plastic variables.  With
\texttt{PLASTIC\_MODEL = F}, the plastic variables are simply irrelevant.

\subsection*{3. Create a cohesive sediment input file}

Copy \texttt{generic\_sediment\_c.inp} or
\texttt{mp\_sediment\_a\_no\_adv.inp} to
\texttt{generic\_sediment\_floc\_only.inp}.  Then make these key edits:

\begin{verbatim}
BEDLOAD = F
SUSLOAD = T
SED_PTSOURCE = T
MORPHO_MODEL = F
MORPHO_FACTOR = 1.0

NSED = 1
SED_NAME = coarse_sand_1
SED_TYPE = cohesive
SED_SD50 = 0.004
SED_WSET = 0.1
SED_SRHO = 2650.
SED_ERAT = 0.58E-3
SED_TAUE = 0.1
SED_TAUD = 0.1
SED_PORS = 0.5

NBED = 1
INF_BED = T
INIT_BED_POROSITY = 0.50
INIT_BED_THICKNESS = 10.0
INIT_BED_FRACTION = 1.0
\end{verbatim}

\texttt{SED\_SD50 = 0.004} mm is a placeholder primary-particle scale
consistent with the floc literature values used in \texttt{mod\_flocmod.F}.
For Delaware/Schuylkill suspended sediment, this should be tuned to the
available grain-size data.  The important code-facing point is that it should
be much smaller than the current sand values \texttt{0.1--1.0} mm.

\section*{What Not To Expect From This Minimum Path}

\begin{itemize}
  \item It will not output an explicit floc diameter or full floc size
        distribution.  The source only stores the floc-affected settling
        velocity internally in \texttt{sed(i)\%wset0}; that field is not
        registered as a standard sediment NetCDF output in the current
        \texttt{mod\_ncdio.F}.
  \item It will not use \texttt{SED\_CSED} for initial water-column sediment.
        The ORIG initialization sets suspended sediment concentration to zero;
        the practical sediment source is the river forcing unless
        \texttt{init\_sed.F} is edited.
  \item In the cohesive settling flux routine, the critical deposition stress
        gate is commented out, so \texttt{SED\_TAUD} should not be relied on to
        suppress deposition under high shear in this path.
  \item The floc enhancement activates only where the simulated salinity is
        between 5 and 15 PSU.  River-source salinity in the forcing file is
        zero, so floc enhancement is expected mainly in the brackish mixing
        region, not at the fresh river nodes.
\end{itemize}

\section*{Why Not Use \texttt{mod\_flocmod.F} First?}

\texttt{mod\_flocmod.F} is the more complete sediment floc model.  It computes
transfers among cohesive sediment classes through shear aggregation,
differential-settling aggregation, shear fragmentation, optional
collision-induced fragmentation, and optional bed deflocculation.  However, it
is only active in the CSTMS build:
\begin{verbatim}
#if defined (SEDIMENT) && (CSTMS_SED)
\end{verbatim}
and is called from \texttt{mod\_sed\_cstms.F} only when \texttt{SED\_FLOCS =
T}.  Using it would require:
\begin{itemize}
  \item rebuilding with \texttt{-DCSTMS\_SED} instead of \texttt{-DORIG\_SED};
  \item writing a CSTMS-format sediment input with \texttt{NCS}, \texttt{NNS},
        \texttt{MUD\_*}, \texttt{SED\_FLOCS}, and \texttt{f\_*} parameters;
  \item creating river NetCDF variables for every cohesive class name;
  \item rechecking NetCDF output and restart behavior for the CSTMS variables.
\end{itemize}

That is a valid future path if the research question requires an evolving floc
size distribution.  It is not the minimum-effort path for a first estuary-scale
sediment-floc-only experiment.

\section*{Recommended First Experiment}

Run a short pair of tests using the same hydrodynamics and river sediment
forcing:
\begin{enumerate}[(1)]
  \item \textbf{No-floc control:} one sediment class with the same
        \texttt{SED\_SD50} and \texttt{SED\_WSET}, but
        \texttt{SED\_TYPE = non-cohesive}.
  \item \textbf{Floc case:} identical file except
        \texttt{SED\_TYPE = cohesive}.
\end{enumerate}

Compare the output variable \texttt{coarse\_sand\_1} in surface, bottom, and
depth-averaged views, plus \texttt{coarse\_sand\_1\_bedfrac}.  This isolates
the effect of the ORIG floc-settling formula without invoking plastic
aggregation or the larger CSTMS framework.

\section*{Bottom Line}

For minimum effort, do not modify source code and do not start with
\texttt{mod\_flocmod.F}.  Use the existing \texttt{ORIG\_SED} module, turn off
plastic, keep the existing river sediment variable name, and run one cohesive
sediment class.  This gives a sediment-only floc sensitivity experiment
immediately.  Use CSTMS later only if the goal changes from ``floc-affected
settling'' to an explicit evolving floc population.

\end{document}
